% Compiled with LuaLaTeX. Main font is Roboto, loaded from the system TTFs via
% fontspec/fontconfig (the roboto LaTeX package is not installed here).
% Layout based on the Rover Resume base template
% (https://github.com/subidit/rover-resume): ruled section headers and a
% strict single-level bullet hierarchy, with run-in project titles that each
% collect their supporting artifacts as inline links.
\documentclass[11pt,a4paper]{article}

\usepackage[a4paper,margin=0.72in]{geometry}
\usepackage{fontspec}
\setmainfont{Roboto}[
  Ligatures=TeX,
  BoldFont={Roboto Medium},
  ItalicFont={Roboto Italic},
  BoldItalicFont={Roboto Medium Italic},
]
% Math is set in Fira Math (loaded from the system OTF), which pairs with the
% geometric-humanist Roboto text. unicode-math must follow fontspec.
\usepackage{unicode-math}
\setmathfont{Fira Math}
\usepackage[dvipsnames]{xcolor}
\definecolor{linkblue}{HTML}{1155CC}
\usepackage[colorlinks=true,urlcolor=linkblue,linkcolor=linkblue]{hyperref}
\hypersetup{
  pdftitle={Nikita Kazeev - CV},
  pdfauthor={Nikita Kazeev},
  pdfsubject={Research Scientist focusing on Reliability and Uncertainty in Agentic Systems},
  pdfkeywords={ML \& AI, Python, PyTorch, C++, HPC, Physics, Research Writing \& Speaking, Leadership \& Mentorship},
}
\usepackage{titlesec}
\usepackage{enumitem}
\usepackage{xstring}        % substring test for the equal-contribution marker

% Publications are stored in BibLaTeX format (src/data/publications.bib) and
% typeset with biber. A numeric style with file-order sorting reproduces the
% manually numbered list the CV used previously.
% url/doi/eprint are suppressed from standalone printing (url=false etc.); the
% title hyperlinks to the entry's url instead (see \DeclareFieldFormat below).
% Custom fields (sessionurl, eqcontrib) are declared in cvdm.dbx, loaded via the
% datamodel option — current biblatex forbids \DeclareDatamodel* in the preamble.
\usepackage[backend=biber,datamodel=cvdm,style=numeric,sorting=none,maxbibnames=99,giveninits=false,doi=false,url=false,isbn=false,eprint=false]{biblatex}
\addbibresource{publications.bib}
\renewcommand*{\bibfont}{\small}
\setlength{\bibitemsep}{3pt}
\setlength{\bibparsep}{0pt}
\defbibheading{none}{}


% Unnumbered publication list: a flush-left hanging list with no [n] labels.
\defbibenvironment{bibliography}
  {\list{}{\setlength{\leftmargin}{1.2em}\setlength{\itemindent}{-\leftmargin}\setlength{\itemsep}{\bibitemsep}\setlength{\parsep}{\bibparsep}}}
  {\endlist}
  {\item}
% Boldface my own name wherever it appears in an author list, and append a
% superscript star to any author named in the entry's eqcontrib field.
\renewcommand*{\mkbibcompletename}[1]{\ifboolexpr{ test {\ifdefstring{\namepartfamily}{Kazeev}} }{\textbf{#1}}{#1}\iffieldundef{eqcontrib}{}{\IfSubStr{\thefield{eqcontrib}}{\namepartfamily}{\textsuperscript{*}}{}}}
\DeclareFieldFormat{addendum}{\textbf{#1}}
% Make the (quoted) title the clickable link to the publication.
\DeclareFieldFormat[article,inproceedings]{title}{\iffieldundef{url}{\mkbibquote{#1\isdot}}{\href{\thefield{url}}{\mkbibquote{#1\isdot}}}}
% ICML 2026 entries: lead with the session time (bold) so a reader sees when
% to attend, drop the redundant standalone "July 2026" date, and remove the
% note from its usual trailing position. Other entries are untouched.
\AtEveryBibitem{\ifkeyword{icml2026}{\clearfield{year}\clearfield{month}\clearfield{day}\printtext{\textbf{\thefield{note}}\addspace\textemdash\addspace}\clearfield{note}}{}}

% Variant filters: the build script bakes the current variant tag into the
% template, so `safety` is a literal keyword here (empty for the
% generic CV, which prints everything).
\defbibfilter{cvicml}{ \keyword{icml2026} and \keyword{safety} }
\defbibfilter{cvselected}{ not \keyword{icml2026} and \keyword{safety} }

\setcounter{secnumdepth}{0}            % no section numbering
\titleformat{\section}{\large\bfseries\scshape}{}{0em}{}[\titlerule]
\titlespacing*{\section}{0pt}{8pt}{4pt}
\setlist[itemize]{topsep=1pt,itemsep=0pt,parsep=0pt,leftmargin=1.4em,label=\textbullet}
\setlength{\parindent}{0pt}
\pagestyle{empty}

% Run-in list of supporting artifacts after a project title.
\newcommand{\artifacts}[1]{\hspace{0.45em}{\footnotesize #1}}

% Projects are indented under their employer to establish visual hierarchy:
% the employer line is the flush-left anchor, projects nest one step in.
\newlength{\projindent}\setlength{\projindent}{1.3em}
\newlength{\projbulletindent}\setlength{\projbulletindent}{2.7em}

\begin{document}

% ===== Header =====
\begin{center}
  {\LARGE\bfseries Nikita Kazeev}\\[3pt]
  {\small
    \href{mailto:kazeevn@gmail.com}{kazeevn@gmail.com} \textperiodcentered{}
    \href{https://kazeevn.github.io}{Website} \textperiodcentered{}
    \href{https://github.com/kazeevn}{GitHub} \textperiodcentered{}
    \href{https://linkedin.com/in/nikita-kazeev}{LinkedIn} \textperiodcentered{}
    \href{https://scholar.google.com/citations?user=vamy2okAAAAJ}{Google Scholar} \textperiodcentered{}
    \href{https://orcid.org/0000-0002-5699-7634}{ORCID}%
  }\\[5pt]
  \textit{ Research Scientist focusing on Reliability and Uncertainty in Agentic Systems }
\end{center}

% ===== PAGE 1: Work Experience & Projects =====
\section{Work Experience}
\textsc{\textbf{\boldmath{}Perfomax}}\quad{\small\color{black!55}\itshape AI Advisor} \hfill\textbf{2026--present}\par
\smallskip

\textsc{\textbf{\boldmath{}Constructor Group}}\quad{\small\color{black!55}\itshape Machine Learning Consultant} \hfill\textbf{2021--present}\par
\smallskip

\textsc{\textbf{\boldmath{}National University of Singapore}}\quad{\small\color{black!55}\itshape Postdoc under \href{https://en.wikipedia.org/wiki/Konstantin\_Novoselov}{Kostya Novoselov}} \hfill\textbf{2022--present}\par
\vspace{2pt}
{\leftskip\projindent\textbf{Leadership}\par}
\begin{itemize}[leftmargin=\projbulletindent]
  \item Program Chair for the \href{https://ai4x.cc/}{AI4X 2025/6 conferences}
  \item Main organiser of the \href{https://multiscale-ai.github.io/}{ICLR 2025 Workshop on Machine Learning Multiscale Processes}
  \item Co-PI of a US\$3.4m AI Singapore grant on multiscale machine learning
\end{itemize}
\vspace{2pt}
{\leftskip\projindent\textbf{LLM Agent Safety \& Alignment}\artifacts{ \href{https://openreview.net/forum?id=zqShTtnTZc}{ICML 2026 Workshop} \textperiodcentered{} \href{https://openreview.net/forum?id=J3Vqv2fFse}{ICML 2026 Workshop} \textperiodcentered{} \href{https://github.com/kazeevn/referenece-checker}{GitHub} }\par}
\begin{itemize}[leftmargin=\projbulletindent]
  \item Conceived and developed a novel geometric framework to predict multi-step reasoning failures in LLM agents using semantic abstraction trajectories
  \item Designed uncertainty quantification (UQ) algorithms for LLM reasoning loops, improving agent error detection and calibration
  \item Developed and open-sourced \href{https://github.com/kazeevn/referenece-checker}{reference-checker}, a tool for verifying claims and citation references in scientific papers to detect and mitigate LLM hallucinations
  \item Authored a position paper on Pluralistic Alignment, advocating for aligning AI models with aspirational human values rather than replicating behavioral flaws
\end{itemize}
\vspace{2pt}
{\leftskip\projindent\textbf{Wyckoff Transformer, a generative model for materials}\artifacts{ \href{https://icml.cc/virtual/2025/poster/44595}{ICML 2025} \textperiodcentered{} \href{https://github.com/SymmetryAdvantage/WyckoffTransformer}{GitHub} }\par}
\begin{itemize}[leftmargin=\projbulletindent]
  \item Alignment by design: enforced physical principles directly at the architectural level as hard constraints, mathematically guaranteeing plausible generations
  \item Solved a systemic failure mode where generative models default to asymmetric outputs despite symmetric training data, mitigating this structural bias via coordinate-free representations
  \item Increased materials discovery yield (S.U.N.) by 1.24x while cutting inference time by 4 orders of magnitude (to 0.05 GPU-ms per structure)
  \item Orchestrated large-scale (10k structures) DFT computation loops to verify out-of-distribution stability; led a team of 6
\end{itemize}
\smallskip

\textsc{\textbf{\boldmath{}$\langle$CERN $\mid$ Yandex $\rightarrow$ HSE University$\rangle$}}\quad{\small\color{black!55}\itshape Researcher} \hfill\textbf{2014--2022}\par
{\small \href{https://breakthroughprize.org/Laureates/1/L3995}{2025 Breakthrough Prize in Fundamental Physics} as a member of LHCb. }\par
\vspace{2pt}
{\leftskip\projindent\textbf{Generative models uncertainty estimation}\artifacts{ \href{https://iopscience.iop.org/article/10.1088/1742-6596/2438/1/012088}{paper} \textperiodcentered{} \href{https://cds.cern.ch/record/2837844}{conference} \textperiodcentered{} \href{https://gitlab.com/lambda-hse/lhcb-rich-gan-uncertainty}{code 1} \textperiodcentered{} \href{https://gitlab.com/lambda-hse/gan-uncertainty-ensembles}{code 2} }\par}
\begin{itemize}[leftmargin=\projbulletindent]
  \item Developed the first methods for estimating uncertainty of conditional GANs
  \item Led a team of 3 students
\end{itemize}
\vspace{2pt}
{\leftskip\projindent\textbf{Machine learning on noisy data}\artifacts{ \href{http://doi.org/10.1088/1748-0221/14/08/P08020}{paper 1} \textperiodcentered{} \href{https://ml4physicalsciences.github.io/2019/files/NeurIPS\_ML4PS\_2019\_122.pdf}{paper 2} \textperiodcentered{} \href{https://github.com/yandexdataschool/ML-sWeights-experiments/}{code} }\par}
\begin{itemize}[leftmargin=\projbulletindent]
  \item Derived a statistically rigorous method to train classifiers on sPlot-weighted (label-noisy) data, standard in high-energy physics
\end{itemize}
\vspace{2pt}
{\leftskip\projindent\textbf{CatBoost: Combatting Biases in Boosting}\artifacts{ \href{https://arxiv.org/abs/1706.09516v1}{paper} }\par}
\begin{itemize}[leftmargin=\projbulletindent]
  \item Designed and built distributed experiment infrastructure to scale systematic bias analysis across large-scale tabular datasets
  \item Investigated dynamic boosting formulations to eliminate target leakage and prediction shift, mitigating statistical bias in gradient boosting
\end{itemize}
\smallskip

\newpage

% ===== PAGE 2: Education & Activities =====
\section{Education}
\textbf{HSE University \& Sapienza Università di Roma}\hfill\textbf{2016--2020}\par
\textit{PhD in Computer Science and Physics (Double Degree)}\par
{\leftskip\projindent Supervisors: \href{https://www.linkedin.com/in/andreyustyuzhanin/}{Andrey Ustyuzhanin} \& \href{https://it.linkedin.com/in/barbara-sciascia-234935a}{Barbara Sciascia}\par}
{\leftskip\projindent Thesis: \href{https://repository.cern/records/034yx-nt191}{Machine Learning for particle identification in the LHCb detector}\par}
\smallskip

\textbf{Product Management course at Yandex}\hfill\textbf{2018}\par
\textit{Focused coursework in product strategy and execution.}\par
\smallskip

\textbf{Yandex School of Data Analysis}\hfill\textbf{2013--2015}\par
\textit{Master's-level CS programme}\par
{\leftskip\projindent Algorithms, machine learning, deep learning, and distributed systems.\par}
\smallskip

\textbf{Moscow Institute of Physics and Technology}\par
\textit{MS in Physics}\hfill\textbf{2014--2016}\par
{\leftskip\projindent Optimisation of data processing of the LHCb experiment.\par}
\textit{BS in Physics}\hfill\textbf{2010--2014}\par
{\leftskip\projindent Wave-packet molecular dynamics of electrons in nonideal plasma.\par}
\smallskip

\textbf{International Junior Science Olympiad (IJSO)}\hfill\textbf{Korea, 2008}\par
{\leftskip\projindent Silver medal\par}
\smallskip


\section{Mentorship}
\begin{itemize}
  \item Mentored 9 students (\href{https://www.hse.ru/edu/vkr/219430466}{1}, \href{https://www.hse.ru/edu/vkr/296294066}{2}, \href{https://www.hse.ru/edu/vkr/219524130}{3}, \href{https://www.hse.ru/en/edu/vkr/470561417}{4}, \href{https://www.hse.ru/en/edu/vkr/470935224}{5}, \href{https://www.hse.ru/ma/datasci/students/diplomas/631565328}{6}, \href{https://www.hse.ru/ma/datasci/students/diplomas/631565422}{7}, \href{https://www.hse.ru/edu/vkr/624890096}{8}, \href{https://www.hse.ru/ba/ami/students/diplomas/834351454}{9}), 3 interns, and student workshop projects (\href{https://sochisirius.ru/news/3148}{1}, \href{https://web.archive.org/web/20220512040654/https://academy.yandex.ru/posts/vossozdat-arkhitekturu-bazy-dannykh-i-nauchit-golosovogo-pomoschnika-predskazyvat-namereniya-polzovatelya}{2}).
\end{itemize}
\section{Teaching \& Outreach}
\begin{itemize}
  \item Pitch research to the general public through major newspaper coverage (\href{https://www.kommersant.ru/doc/7832307}{1}, \href{https://www.kommersant.ru/doc/6122817}{2}), blog posts (\href{https://ifim.nus.edu.sg/tailor-made-2d-materials-became-closer-with-the-new-machine-learning-algorithm-developed-at-nus/}{1}, \href{https://communities.springernature.com/posts/sparse-representation-for-machine-learning-the-properties-of-defects-in-2d-materials}{2}), a \href{https://elementy.ru/events/435506/IV\_Nauchnye\_boi\_NIU\_VShE\_na\_Dne\_Vyshki}{science slam}, interviews (\href{https://web.archive.org/web/20210518043835/https://academy.yandex.ru/posts/beresh-i-razbiraeshsya}{1}, \href{https://habr.com/en/company/yandex/blog/422761/}{2}), and a public \href{https://www.youtube.com/watch?v=zLFZ\_z-7-JI&feature=youtu.be&t=7779}{talk}.
  \item Lecturer in the \href{https://www.hse.ru/en/news/edu/259378493.html}{award-winning} Coursera Advanced Machine Learning \href{https://web.archive.org/web/20220127183223/https://www.coursera.org/specializations/aml}{  specialisation}
  \item Delivered more than 20 \href{https://docs.google.com/document/d/1IOlBse4q0YgEoL\_RCJTvJfbFhFaD\_YMAIJJ3W77B\_X0/edit?usp=sharing}{conference talks}.
\end{itemize}
\section{Service}
\begin{itemize}
  \item Reviewer for NeurIPS, RSC Advances, Machine Learning: Science and Technology, Digital Discovery
  \item Peer Staff Supporter at NUS, serving as a first-line support for mental wellbeing.
\end{itemize}

\newpage

% ===== PAGE 3: Skills & Publications =====
\section{Skills}
\begin{itemize}
  \item \textbf{Agent Reliability \& UQ} --- developed geometric and statistical frameworks to predict LLM agent reasoning failures using semantic abstraction trajectories.
  \item \textbf{Aspirational Alignment \& Verification} --- formulated alignment strategies based on high-level human values and established benchmarks (e.g. Foresight-Phys) for validating AI systems in complex environments.
  \item \textbf{Uncertainty quantification} --- built the first methods for estimating the uncertainty of conditional GANs; calibrated generative and discriminative models.
  \item \textbf{Inference and systems at scale} --- C++ production integration, large-scale DFT/VASP on HPC clusters, and reproducible experiment pipelines.
  \item \textbf{Research leadership} --- lead multidisciplinary teams, chair conferences, and mentor students from idea to publication.
\end{itemize}
{\small\centering ML \& AI \textbullet\ Python \textbullet\ PyTorch \textbullet\ C++ \textbullet\ HPC \textbullet\ Physics \textbullet\ Research Writing \& Speaking \textbullet\ Leadership \& Mentorship\par}

% Pull every entry from publications.bib into the bibliography; the filters
% below split them into the two sections.
\nocite{*}
\section{Presenting at ICML 2026}
\printbibliography[filter=cvicml,heading=none,resetnumbers=true]

\section{Selected Publications}
\printbibliography[filter=cvselected,heading=none,resetnumbers=true]

\end{document}